\section{Problem Definition}
\label{sec:problem}

\subsection{The Communication Gap}

Executives, product managers, consultants, and program leads communicate plans through visual timelines. These artifacts appear in board presentations, investor decks, quarterly reviews, transformation roadmaps, and architecture evolution documents. The quality of these visuals directly affects stakeholder comprehension and decision-making.

Yet producing presentation-quality timelines remains surprisingly difficult. The gap exists because available tools fall into two categories, each failing half the requirement:

\begin{description}
  \item[Documentation Tools] (Mermaid, PlantUML, text-based DSLs) — machine-readable, version-controllable, agent-friendly, but produce diagrams suited for technical documentation rather than executive communication.
  
  \item[Design Tools] (PowerPoint, Visio, Lucidchart, think-cell, Figma) — produce polished visuals, but require significant manual effort, resist automation, and cannot be reliably generated by AI agents.
\end{description}

Timeline Compiler bridges this gap: a compiler target for structured plan data that produces presentation-quality visual output.

\subsection{Why Existing Tools Are Insufficient}

\subsubsection{Mermaid and Documentation-Grade Diagramming}

Mermaid~\cite{mermaid2023} has become the de facto standard for diagrams-as-code in technical documentation. Its Gantt syntax illustrates the limitation:

\begin{lstlisting}[language=yaml,caption={Mermaid Gantt — functional but not presentation-grade}]
gantt
    title Product Roadmap Q3
    dateFormat YYYY-MM-DD
    section Platform
    Core API :a1, 2026-07-01, 30d
    Mobile SDK :a2, after a1, 45d
    section Integrations  
    Partner A :b1, 2026-08-01, 60d
\end{lstlisting}

Mermaid excels at documentation — syntax is simple, output is deterministic, rendering integrates with GitHub, GitLab, and Notion. But the output is not presentation-grade:

\begin{itemize}
  \item Fixed visual style optimized for legibility in documentation contexts
  \item No theming beyond basic color overrides
  \item Limited typographic control
  \item Dependency-centric (task scheduling) rather than communication-centric (timeline narrative)
  \item No built-in support for swimlanes as first-class visual groupings
  \item No annotations, callouts, or emphasis mechanisms for executive communication
\end{itemize}

The result: Mermaid diagrams are useful in READMEs and wikis but inappropriate for board decks and investor presentations.

\subsubsection{Project Management Tools}

Tools like Microsoft Project, Jira, Asana, Monday.com, and Linear are optimized for \emph{managing work}, not \emph{communicating plans}. Their timeline views:

\begin{itemize}
  \item Export poorly or not at all to static visual formats
  \item Embed project-management semantics (resources, dependencies, critical paths) that clutter executive communication
  \item Assume ongoing interaction rather than snapshot communication
  \item Produce visuals tied to the tool's aesthetic, not the organization's brand
\end{itemize}

A quarterly roadmap exported from Jira looks like a Jira export, not like a McKinsey slide.

\subsubsection{Design and Presentation Tools}

PowerPoint, Keynote, Visio, Figma, and Lucidchart produce beautiful results — but at the cost of manual effort. Each timeline is a bespoke artifact:

\begin{itemize}
  \item No structured data input; changes require manual redrawing
  \item Not version-controllable in meaningful ways (binary or JSON blobs)
  \item Not agent-accessible; an AI cannot reliably produce a Visio file
  \item Theming requires manual application
  \item Determinism is impossible; two designers will produce different artifacts from the same brief
\end{itemize}

This creates a maintenance burden. When scope shifts, the timeline must be manually updated. When quarterly plans roll forward, someone rebuilds the slide.

\subsubsection{AI Agent Blindspot}

Modern AI agents (LLMs with tool use, planning agents, copilots) can generate structured plans — sprints, phases, milestones, dependencies. But they have no widely-adopted open-source target for rendering those plans into visual artifacts.

An agent can output JSON or YAML describing a roadmap. But then what? Without a compiler that accepts structured input and produces presentation-grade output, the agent's plan stops at structured text.

\subsection{Who This Is For}

Timeline Compiler serves two complementary audiences:

\subsubsection{Human Authors}

\begin{description}
  \item[Executives and Chiefs of Staff] — Need polished quarterly plans, transformation roadmaps, and board-ready timelines without designer dependency.
  
  \item[Product Managers] — Need roadmaps that communicate priority and sequencing to stakeholders, not dependency graphs for engineers.
  
  \item[Consultants] — Need branded deliverables (BCG-style roadmaps, McKinsey program timelines) that are data-driven yet visually refined.
  
  \item[DevRel and Technical Writers] — Need version-controlled timelines that embed in documentation but also export to slides.
  
  \item[Program Managers] — Need milestone views and phase summaries that roll up operational detail into executive narrative.
\end{description}

\subsubsection{AI Agents}

AI agents require:
\begin{itemize}
  \item A well-defined intermediate representation (IR) they can generate
  \item Deterministic rendering — the same IR always produces the same visual
  \item Theming separation — agents produce structure, not style
  \item Validation — agents can verify their output conforms to the schema
  \item Embeddability — agents can invoke rendering as a tool
\end{itemize}

Timeline Compiler's IR is designed as a \emph{compiler target}: a stable, documented, validatable representation that agents can emit and humans can inspect, version, and override.

\subsection{What Success Looks Like}

\subsubsection{Scenario: Product Roadmap}

A product manager writes (or an agent generates):

\begin{lstlisting}[language=yaml,caption={Product Roadmap IR Example}]
timeline:
  title: "Platform Roadmap 2026"
  range: { start: 2026-Q1, end: 2026-Q4 }
  tracks:
    - id: platform
      label: "Core Platform"
      activities:
        - label: "API v2"
          range: { start: 2026-01, end: 2026-03 }
          status: complete
        - label: "SDK Release"
          range: { start: 2026-04, end: 2026-06 }
          status: in-progress
          progress: 0.6
    - id: mobile
      label: "Mobile"
      activities:
        - label: "iOS App"
          range: { start: 2026-05, end: 2026-08 }
          status: planned
  milestones:
    - label: "Public Beta"
      date: 2026-06-15
      tracks: [platform, mobile]
\end{lstlisting}

The compiler, given a theme (e.g., \texttt{corporate-blue}), produces a PDF suitable for a board presentation: clean typography, proper spacing, consistent color palette, no visual noise.

\subsubsection{Scenario: Transformation Plan}

A management consultant defines a multi-year digital transformation:

\begin{lstlisting}[language=yaml,caption={Transformation Timeline IR Example}]
timeline:
  title: "Digital Transformation — 3 Year Plan"
  range: { start: 2026-01, end: 2028-12 }
  sections:
    - id: foundation
      label: "Foundation"
      range: { start: 2026-01, end: 2026-12 }
    - id: scaling
      label: "Scaling"  
      range: { start: 2027-01, end: 2027-12 }
    - id: optimization
      label: "Optimization"
      range: { start: 2028-01, end: 2028-12 }
  tracks:
    - id: technology
      label: "Technology"
      activities:
        - label: "Cloud Migration"
          range: { start: 2026-03, end: 2026-09 }
          section: foundation
        - label: "Platform Modernization"
          range: { start: 2027-01, end: 2027-09 }
          section: scaling
    - id: organization
      label: "Organization"
      activities:
        - label: "Team Restructuring"
          range: { start: 2026-06, end: 2026-12 }
          section: foundation
  annotations:
    - label: "Board Checkpoint"
      date: 2026-12-15
      style: callout
\end{lstlisting}

The output: a three-year view with clear phase demarcation, swimlanes for workstreams, and a visual hierarchy that communicates strategic sequencing rather than tactical dependencies.

\subsubsection{Success Criteria}

Timeline Compiler succeeds when:

\begin{enumerate}
  \item A human can write the IR by hand in under 10 minutes for a typical roadmap
  \item An AI agent can generate valid IR without special prompting beyond the schema
  \item The same IR, rendered with the same theme, produces byte-identical output across runs (determinism)
  \item A non-designer can produce board-ready visuals by selecting a theme
  \item The IR is diff-friendly — changes to the plan produce readable diffs
  \item The output is accepted in contexts where PowerPoint slides are currently required
\end{enumerate}

\subsection{What This Is Not}

To prevent scope creep and misunderstanding:

\begin{itemize}
  \item \textbf{Not a project management tool} — does not track work, assign resources, or compute schedules
  \item \textbf{Not a Gantt chart replacement} — though it can render timeline views, it is not optimized for dependency visualization
  \item \textbf{Not a Jira/Asana competitor} — does not replace work-tracking systems; may consume data from them
  \item \textbf{Not a design tool} — does not provide a canvas, drag-and-drop, or WYSIWYG editing
  \item \textbf{Not interactive} — output is static (SVG, PNG, PDF, PPTX); interactivity is out of scope
\end{itemize}

Timeline Compiler is a \emph{compiler}: structured input in, presentation-quality visual output out. The value is in the deterministic, themeable, agent-accessible transformation — not in the tooling around authoring or managing plans.
